\documentclass[12pt]{article}
\usepackage{inputenc}
\usepackage[spanish]{babel}
\decimalpoint % Usar puntos (.) en lugar de comas (,) como separadores decimales
\usepackage{geometry}
\usepackage{amsmath} % Para entornos matemáticos avanzados
\usepackage{amssymb} % Para el uso de \mathbb
\usepackage{amsthm} % Para entornos de teoremas
\usepackage{hyperref} % Para hipervínculos

\geometry{margin=1in}

% Definición de entornos personalizados
\newtheorem{problem}{Problema}

\title{Procesos Estocásticos I \\ Tarea 1}
\author{Semestre 2026-2}
\date{}

\begin{document}

\maketitle

\textbf{Instrucciones:} Resuelva los siguientes problemas. Justifique sus respuestas claramente y muestre todos los pasos. Puede utilizar los conceptos y resultados presentados en las notas de clase.

\begin{problem}
\textbf{Definición y propiedades de cadenas de Markov.}

Considere una cadena de Markov homogénea $\{X_n : n \in \mathbb{N}\}$ con espacio de estados $S = \{1, 2, 3\}$ y matriz de transición:
\[
P = \begin{pmatrix}
0.4 & 0.3 & 0.3 \\
0.2 & 0.5 & 0.3 \\
0.1 & 0.4 & 0.5
\end{pmatrix}.
\]

\begin{enumerate}
    \item[(a)] Verifique que $P$ es una matriz estocástica. ¿Cuáles son las condiciones que debe satisfacer?
    
    \item[(b)] Si la distribución inicial es $\pi = (0.3, 0.4, 0.3)$, calcule la distribución de probabilidades en el tiempo $n=1$.
    
    \item[(c)] Interprete el valor $\pi_1^{(1)}$ en términos de probabilidades en la cadena.
\end{enumerate}
\end{problem}

\begin{problem}
\textbf{Ecuación de Chapman-Kolmogorov y probabilidades en $n$ pasos.}

Nuevamente considere la cadena de Markov del Problema 1 con matriz de transición:
\[
P = \begin{pmatrix}
0.4 & 0.3 & 0.3 \\
0.2 & 0.5 & 0.3 \\
0.1 & 0.4 & 0.5
\end{pmatrix}.
\]

\begin{enumerate}
    \item[(a)] Calcule $P^2 = P \times P$ y encuentre el elemento $(P^2)_{2,1}$. ¿Qué representa este valor?
    
    \item[(b)] Utilice la ecuación de Chapman-Kolmogorov para calcular $p_{1,3}^{(2)}$ (la probabilidad de ir del estado 1 al estado 3 en exactamente 2 pasos). Verifique su resultado con el elemento $(P^2)_{1,3}$.
    
    \item[(c)] Enuncie formalmente la ecuación de Chapman-Kolmogorov y explique su interpretación intuitiva.
\end{enumerate}
\end{problem}

\begin{problem}
\textbf{Caminata aleatoria y su propiedad de Markov.}

Sea $\{Y_i : i \in \mathbb{N}\}$ una sucesión de variables aleatorias independientes e idénticamente distribuidas con distribución $\mathbb{P}(Y_i = 1) = p$ y $\mathbb{P}(Y_i = -1) = 1-p$, donde $0 < p < 1$.

Defina la caminata aleatoria como:
\[
X_n = \sum_{i=1}^n Y_i, \quad \text{con } X_0 = 0.
\]

\begin{enumerate}
    \item[(a)] Demuestre que $\{X_n\}$ satisface la propiedad de Markov, es decir, que es una cadena de Markov homogénea.
    
    \item[(b)] Calcule las probabilidades de transición $p_{i,i+1}$ y $p_{i,i-1}$ para cualquier estado $i$.
    
    \item[(c)] Suponga que $p = 0.6$. Calcule $\mathbb{P}(X_2 = 0)$ y $\mathbb{P}(X_2 = 2)$ partiendo de $X_0 = 0$.
    
    \item[(d)] ¿Por qué es importante que los incrementos $Y_i$ sean independientes para que $\{X_n\}$ sea una cadena de Markov?
\end{enumerate}
\end{problem}

\begin{problem}
\textbf{Diagonalización de la matriz de transición y comportamiento asintótico.}

Considere una cadena de Markov de dos estados con matriz de transición:
\[
P = \begin{pmatrix}
1 - \alpha & \alpha \\
\beta & 1 - \beta
\end{pmatrix},
\]
donde $\alpha = 0.2$ y $\beta = 0.3$.

\begin{enumerate}
    \item[(a)] Calcule los valores propios de $P$. (Sugerencia: Uno de ellos es $\lambda_1 = 1$.)
    
    \item[(b)] Encuentre los vectores propios asociados a cada valor propio.
    
    \item[(c)] Escriba la factorización $P = VDV^{-1}$ donde $V$ es la matriz de vectores propios y $D$ es la matriz diagonal de valores propios.
    
    \item[(d)] Utilizando la factorización anterior, calcule $P^n$ de forma explícita y determine $\lim_{n \to \infty} P^n$.
    
    \item[(e)] Interprete el significado de $\lim_{n \to \infty} P^n$. ¿Qué nos dice sobre el comportamiento a largo plazo de la cadena?
    
    \item[(f)] Identifique la distribución estacionaria $\pi^*$ de la cadena (es decir, la distribución tal que $\pi^* P = \pi^*$).
\end{enumerate}
\end{problem}

\begin{problem}
\textbf{Propiedad de Markov: independencia del futuro del pasado.}

Sea $\{X_n : n \in \mathbb{N}\}$ una cadena de Markov homogénea con espacio de estados $S$. Sea $i, j, k \in S$ estados arbitrarios y sean $m, n \in \mathbb{N}$ con $m < n$.

\begin{enumerate}
    \item[(a)] Demuestre que:
    \[
    \mathbb{P}(X_{n} = j \mid X_{n-1} = i, X_{m} = k) = \mathbb{P}(X_{n} = j \mid X_{n-1} = i).
    \]
    
    Interprete este resultado: ¿qué significa que el lado izquierdo sea igual al lado derecho?
    
    \item[(b)] Generalice el resultado anterior y demuestre que para cualquier $n > m$:
    \[
    \mathbb{P}(X_{n} = j \mid X_{n-1} = i, X_{m} = k) = \mathbb{P}(X_{n} = j \mid X_{n-1} = i).
    \]
    
    \item[(c)] Explique cómo este resultado refleja la ``propiedad de falta de memoria'' (``memoryless property'') de las cadenas de Markov.
    
    \item[(d)] ¿Por qué esta propiedad es fundamental para el análisis de cadenas de Markov? Dé un ejemplo de un proceso estocástico que NO satisfaga esta propiedad.
\end{enumerate}
\end{problem}

\begin{problem}
\textbf{Factorización de probabilidades conjuntas en cadenas de Markov.}

Sea $\{X_n : n \in \mathbb{N}\}$ una cadena de Markov homogénea con espacio de estados $S$, distribución inicial $\pi = (\pi_i)_{i \in S}$ y matriz de transición $P = (p_{ij})$.

Considere una secuencia de estados $i_0, i_1, i_2, \ldots, i_n \in S$.

\begin{enumerate}
    \item[(a)] Demuestre que la probabilidad conjunta de observar la secuencia de estados se puede factorizar como:
    \[
    \mathbb{P}(X_n = i_n, X_{n-1} = i_{n-1}, \ldots, X_1 = i_1, X_0 = i_0) = \pi_{i_0} \cdot p_{i_0, i_1} \cdot p_{i_1, i_2} \cdot \ldots \cdot p_{i_{n-1}, i_n}.
    \]
    
    O de forma compacta:
    \[
    \mathbb{P}(X_0 = i_0, X_1 = i_1, \ldots, X_n = i_n) = \pi_{i_0} \prod_{j=0}^{n-1} p_{i_j, i_{j+1}}.
    \]
    
    Proporcione dos demostraciones: una usando la definición de probabilidad condicional, y otra usando inducción matemática.
    
    \item[(b)] Interprete esta factorización: ¿qué significa en términos de cómo se genera una trayectoria de la cadena?
    
    \item[(c)] Considere la cadena de Markov del Problema 1 (con matriz de transición $P$ y distribución inicial $\pi = (0.3, 0.4, 0.3)$). Calcule la probabilidad de la secuencia específica de estados: $X_0 = 1, X_1 = 2, X_2 = 3$ utilizando la factorización anterior.
    
    \item[(d)] Explique por qué el hecho de que la probabilidad conjunta se factorice de esta manera es una consecuencia directa de la propiedad de Markov.
\end{enumerate}
\end{problem}

\end{document}
